\chapter{Troubleshooting and Support}
\label{ch:troubleshooting}

\section{\TL Website}\label{sec:website}
The website for the \TL packages is located at
\url{http://code.google.com/p/tufte-latex/}.  On our website, you'll find
links to our \smallcaps{svn} repository, mailing lists, bug tracker, and documentation.

\section{\TL Mailing Lists}\label{sec:mailing-lists}
There are two mailing lists for the \TL project:

\paragraph{Discussion list}
The \texttt{tufte-latex} discussion list is for asking questions, getting
assistance with problems, and help with troubleshooting.  Release announcements
are also posted to this list.  You can subscribe to the \texttt{tufte-latex}
discussion list at \url{http://groups.google.com/group/tufte-latex}.

\paragraph{Commits list}
The \texttt{tufte-latex-commits} list is a read-only mailing list.  A message
is sent to the list any time the \TL code has been updated.  If you'd like to
keep up with the latest code developments, you may subscribe to this list.  You
can subscribe to the \texttt{tufte-latex-commits} mailing list at
\url{http://groups.google.com/group/tufte-latex-commits}.

\section{Getting Help}\label{sec:getting-help}
If you've encountered a problem with one of the \TL document classes, have a
question, or would like to report a bug, please send an email to our
mailing list or visit our website.

To help us troubleshoot the problem more quickly, please try to compile your
document using the \docclsopt{debug} class option and send the generated
\texttt{.log} file to the mailing list with a brief description of the problem.



\section{Errors, Warnings, and Informational Messages}\label{sec:tl-messages}
The following is a list of all of the errors, warnings, and other messages generated by the \TL classes and a brief description of their meanings.
\index{error messages}\index{warning messages}\index{debug messages}

% Errors
\docmsg{Error: \doccmd{subparagraph} is undefined by this class.}{%
	The \doccmd{subparagraph} command is not defined in the \TL document classes.
	If you'd like to use the \doccmd{subparagraph} command, you'll need to redefine
	it yourself.  See the ``Headings'' section on page~\pageref{sec:headings} for a
	description of the heading styles availaboe in the \TL document classes.}

\docmsg{Error: \doccmd{subsubsection} is undefined by this class.}{%
	The \doccmd{subsubsection} command is not defined in the \TL document classes.
	If you'd like to use the \doccmd{subsubsection} command, you'll need to
	redefine it yourself.  See the ``Headings'' section on
	page~\pageref{sec:headings} for a description of the heading styles availaboe
	in the \TL document classes.}

\docmsg{Error: You may only call \doccmd{morefloats} twice. See the\par\noindent\ \ \ \ \ \ \ \ Tufte-LaTeX documentation for other workarounds.}{%
	\LaTeX{} allocates 18 slots for storing floats.  The first time
	\doccmd{morefloats} is called, it allocates an additional 34 slots.  The second
	time \doccmd{morefloats} is called, it allocates another 26 slots.
	
	The \doccmd{morefloats} command may only be called two times.  Calling it a
	third time will generate this error message.  See
	page~\pageref{err:too-many-floats} for more information.}

% Warnings
\docmsg{Warning: Option `\docopt{class option}' is not supported -{}- ignoring option.}{%
	This warning appears when you've tried to use \docopt{class option} with a \TL
	document class, but \docopt{class option} isn't supported by the \TL document
	class.  In this situation, \docopt{class option} is ignored.}

% Info / Debug messages
\docmsg{Info: The `\docclsopt{symmetric}' option implies `\docclsopt{twoside}'}{%
	You specified the \docclsopt{symmetric} document class option.  This option automatically forces the \docclsopt{twoside} option as well.  See page~\pageref{clsopt:symmetric} for more information on the \docclsopt{symmetric} class option.}


\section{Package Dependencies}\label{sec:dependencies}
The following is a list of packages that the \TL document
classes rely upon.  Packages marked with an asterisk are optional.
\begin{multicols}{2}
	\begin{itemize}
		\item xifthen
		\item ifpdf*
		\item ifxetex*
		\item hyperref
		\item geometry
		\item ragged2e
		\item chngpage \emph{or} changepage
		\item paralist
		\item textcase
		\item soul*
		\item letterspace*
		\item setspace
		\item natbib \emph{and} bibentry
		\item optparams
		\item placeins
		\item mathpazo*
		\item helvet*
		\item fontenc
		\item beramono*
		\item fancyhdr
		\item xcolor
		\item textcomp
		\item titlesec
		\item titletoc
	\end{itemize}
\end{multicols}
